\documentclass[10pt]{article}
\usepackage[T1]{fontenc}
\usepackage[utf8]{inputenc}
\usepackage[french]{babel}
\usepackage{amsmath,amssymb,mathrsfs,mathpazo}
\usepackage{geometry}
\geometry{a4paper,margin=2cm}
\usepackage{graphicx,array,tabularx,tabularray,colortbl,multirow,multicol,booktabs,enumitem,pifont,ulem,fancybox,diagbox,cancel}
\usepackage{pstricks,pstricks-add,pst-plot,pst-tree,pst-node,pst-text,pst-func,pst-math,pst-bspline,pst-3dplot,pst-eucl,pst-solides3d}
\let\pscalculate\undefined
\usepackage{pst-bezier}
\usepackage{tikz,pgfplots,tkz-tab}
\usetikzlibrary{arrows,arrows.meta,patterns,calc,positioning,trees,decorations.markings,shapes.geometric,shapes.misc,shapes.arrows,intersections,decorations.pathmorphing,decorations.pathreplacing,backgrounds,shadows,babel}
\pgfplotsset{compat=1.18}
\usepgfplotslibrary{fillbetween}
\usepackage[np]{numprint}
\usepackage{listings,makecell,pas-tableur,twoopt}
\let\barre\undefined
\usepackage[pstricks]{tablvar}
\definecolor{darkgreen}{RGB}{0,100,0}
\providecommand{\red}{\color{red}}
\providecommand{\blue}{\color{blue}}
\providecommand{\green}{\color{green}}
\providecommand{\black}{\color{black}}
\providecommand{\white}{\color{white}}
\providecommand{\gray}{\color{gray}}
\providecommand{\cyan}{\color{cyan}}
\providecommand{\magenta}{\color{magenta}}
\providecommand{\yellow}{\color{yellow}}
\providecommand{\violet}{\color{violet}}
\providecommand{\orange}{\color{orange}}
\providecommand{\euro}{\texteuro}
\providecommand{\EUR}{\texteuro}
\providecommand{\ieme}{\textsuperscript{e}}
\providecommand{\degres}{\ensuremath{^\circ}}
\providecommand{\degree}{\ensuremath{^\circ}}
\providecommand{\vv}[1]{\overrightarrow{#1}}
\providecommand{\Coord}[1]{#1}
\providecommand{\decosix}{}
\providecommand{\starredbullet}{\ensuremath{\bullet}}
\pagestyle{empty}
\setlength{\parindent}{0pt}
\usepackage{adjustbox}\DeclareUnicodeCharacter{2236}{:}\DeclareUnicodeCharacter{2190}{\ensuremath{\leftarrow}}
\providecommand{\R}{\mathbb{R}}
\providecommand{\N}{\mathbb{N}}
\providecommand{\Z}{\mathbb{Z}}
\providecommand{\Q}{\mathbb{Q}}
\providecommand{\D}{\mathbb{D}}
\providecommand{\C}{\mathbb{C}}
\providecommand{\e}{{\mathrm{e}}}
\providecommand{\vect}[1]{\overrightarrow{#1}}
\providecommand{\vectt}[1]{\overrightarrow{\mathrm{#1}}}
\providecommand{\vvt}[1]{\overrightarrow{\mathrm{#1}}}
\providecommand{\np}[2][]{{#2~#1}}
\providecommand{\vv}[1]{{\overrightarrow{#1}}}
\renewcommand{\d}{{\mathrm{d}}}
\providecommand{\degres}{{^\circ}}
\providecommand{\degre}{{^\circ}}
\providecommand{\blue}{\color{blue}}
\providecommand{\red}{\color{red}}
\providecommand{\enskip}{\quad}
\providecommand{\llbracket}{\lbrack\!\lbrack}
\providecommand{\rrbracket}{\rbrack\!\rbrack}
\providecommand{\euro}{{\text{€}}}
\providecommand{\no}{{\text{n°}}}
\providecommand{\interval}[2]{{[#1;#2]}}
\providecommand{\intervFO}[2]{{[#1;#2[}}
\providecommand{\dotfill}{\ldots}

\providecommand{\ds}{\displaystyle}
\providecommand{\textding}[1]{#1}
\providecommand{\Oij}{\(\mathrm{O};\vec{\imath},\vec{\jmath}\)}
\providecommand{\Oijk}{\(\mathrm{O};\vec{\imath},\vec{\jmath},\vec{k}\)}
\providecommand{\Ouv}{\(\mathrm{O};\vec{u},\vec{v}\)}
\providecommand{\intervFF}[2]{\ensuremath{[#1;#2]}}
\renewcommand{\euro}{\ifmmode\text{\texteuro}\else\texteuro\fi}
\DeclareUnicodeCharacter{00B0}{\ifmmode^{\circ}\else\textdegree\fi}
\DeclareUnicodeCharacter{20AC}{\ifmmode\text{\texteuro}\else\texteuro\fi}
% Énoncés : Éducation nationale. Transcription et corrigés : APMEP, Denis Vergès et contributeurs. Découpage et classement : CoopMaths, Sébastien Lozano, Éric Elter et contributeurs. Adaptation de la présentation : Maths Explorer.
% CC BY-SA : https://coopmaths.fr/www/about/
% https://coopmaths.fr/alea/static/bac/2022/tex/bac_2022_05_sujet1_polynesie_4_cor.tex
\begin{document}
\noindent\textit{APMEP / CoopMaths. CC BY-SA. Présentation adaptée par Maths Explorer.}\par\medskip
\medskip

L'espace est rapporté un repère orthonormal où l'on considère :

\setlength\parindent{1cm}
\begin{itemize}
\item[$\bullet~~$]les points A$(2~;~-1~;~0)$ B$(1~;~0~;~- 3)$, C$(6~;~6~;~1)$  et E$(1~;~2~;~4)$ ;
\item[$\bullet~~$]Le plan $\mathcal{P}$ d'équation cartésienne $2x - y - z + 4 = 0$.
\end{itemize}
\setlength\parindent{0cm}

\medskip

\begin{enumerate}
\item
	\begin{enumerate}
		\item On veut démontrer que le triangle ABC est rectangle en A; on a 
$\vect{\text{AB}}\begin{pmatrix}-1\\1\\-3 \end{pmatrix}$ et $\vect{\text{AC}}\begin{pmatrix}4\\7\\1 \end{pmatrix}$.
		
$\vect{\text{AB}} \cdot \vect{\text{AC}}=-4 +7-3=0$ donc   $\vect{\text{AB}}$ et $\vect{\text{AC}}$ sont orthogonaux et le triangle ABC est donc rectangle en A.

		\item $\vect{\text{BC}}\begin{pmatrix}5\\6\\4 \end{pmatrix}$
donc $\vect{\text{BA}} \cdot \vect{\text{BC}}=5-6+12=11$.

$\text{BA}= \sqrt{1^2+(-1)^2+3^2}=\sqrt{11}$ et $\text{BC}=\sqrt{5^2+6^2+4^2} =\sqrt{77}$

		\item On cherche la mesure en degrés de l'angle $\widehat{\text{ABC}}$ arrondie au degré.

On a alors
\begin{align*}
\vect{\text{BA}} \cdot \vect{\text{BC}}= \text{BA}\times\text{BC}\times \cos\left(\widehat{\text{ABC}}\right)&\iff 11=\sqrt{11}\times \sqrt{77} \times \cos\left(\widehat{\text{ABC}}\right)\\
&\iff 11=\sqrt{11}\times \sqrt{11}\times \sqrt{7}\times \cos\left(\widehat{\text{ABC}}\right) \\
&\iff \cos\left(\widehat{\text{ABC}}\right) = \dfrac{1}{\sqrt{7}}
\end{align*}

La calculatrice donne $\widehat{\text{ABC}} \approx 68\degres$ au degré près.
	\end{enumerate}
	
\item
	\begin{enumerate}
		\item On veut démontrer que le plan $\mathcal{P}$ est parallèle au plan (ABC).

Un vecteur normal du plan $\mathscr{P}$ est le vecteur $\vect{n}\begin{pmatrix}2\\-1\\-1\end{pmatrix}$.

Les vecteurs $\vect{\text{AB}}$ et $\vect{\text{AC}}$ ne sont pas colinéaires car ABC est un triangle rectangle.

De plus $\vect{n} \cdot \vect{\text{AB}}=2\times (-1) -1\times 1 - 1\times (-3) =-2-1+3 = 0$

et

$\vect{n} \cdot \vect{\text{AC}}=2\times 4-1\times 7-1\times 1 = 8 - 7 - 1 = 0$

$\vect{n}$ est donc un vecteur normal à deux vecteurs non colinéaires du plan (ABC), c'est donc un vecteur normal de ce plan et du plan $\mathscr{P}$.

Les plans $\mathscr{P}$ et (ABC) sont donc parallèles.

		\item On déduit une équation cartésienne du plan (ABC).

Le plan (ABC) a donc une équation de la forme $2x - y - z + d = 0$, comme A appartient à ce plan,  on a: $5 + d = 0$ soit $d = -5$.
		
Une équation de (ABC) est donc $2x - y - z - 5 = 0$.

		\item On détermine une représentation paramétrique de la droite $\mathcal{D}$ orthogonale au plan (ABC) et passant par le point E.
		
Un vecteur directeur de la droite $\mathcal{D}$ est donc le vecteur $\vect{n}$ 
		

On a donc $M(x~;~y~;~z) \in \mathcal{D} \iff \vect{\text{E}M} = t \vect{n}$, avec $t \in \R$, soit

$\left\{\begin{array}{l c l}
x - 1&=&2t\\
y - 2&=&- t,\\
z - 4&=&-t
\end{array}\right. t \in \R \iff \left\{\begin{array}{l c l}
x &=& 1 +2t\\
y &=& 2 - t,\\
z &=&4 -t
\end{array}\right. t \in \R$.

		\item On démontre que le projeté orthogonal H  du point E sur le plan (ABC) a pour coordonnées $\left(4~;~\dfrac{1}{2}~;~\dfrac{5}{2}\right)$.
		
		H correspond à l'intersection du plan (ABC) avec la droite perpendiculaire à (ABC) qui passe par E soit la droite $\mathcal{D}$, les coordonnées de H seront donc solutions du système suivant  :
		
\begin{align*}
\begin{cases} x=1+2t\\y=2-t\\z=4-t\\2x-y-z-5=0\end{cases}& \iff\begin{cases} x=1+2t\\y=2-t\\z=4-t\\2(1+2t)-(2-t)-(4-t)-5=0\end{cases} \\
&\iff \begin{cases} x=1+2t\\y=2-t\\z=4-t\\6t-9=0\end{cases} 
&\iff \begin{cases} t=\frac{3}{2}\\x=4\\y = \frac{1}{2}\\z=\frac{5}{2}\end{cases}
\end{align*}
	\end{enumerate}
\item On rappelle que le volume d'une pyramide est donné par $\mathcal{V} = \dfrac13 \mathcal{B}h$ où $\mathcal{B}$ désigne l'aire d'une base et $h$  la hauteur de la pyramide associée à cette base.

On va calculer l'aire du triangle ABC puis démontrer que le volume de la pyramide  ABCE est égal à $16,5$ unités de volume.

$\text{AC}=\sqrt{4^2+7^2+1^2}=\sqrt{66}$ et donc $\mathscr{B}=\dfrac{\text{AB}\times \text{AC}}{2} = \dfrac{\sqrt{11}\times \sqrt{66}}{2}=\dfrac{11\sqrt{6}}{2}$

HE est la hauteur de la pyramide et $\vect{\text{HE}}\begin{pmatrix}-3\\\frac{3}{2}\\\frac{3}{2} \end{pmatrix}$

On a par suite HE $=\sqrt{(-3)^2+\left(\dfrac{3}{2}\right)^2+\left(\dfrac{3}{2}\right)^2}=\sqrt{\dfrac{27}{2}}$ puis

$ \mathscr{V}=\dfrac{1}{3}\times \dfrac{11\sqrt{6}}{2}\times \sqrt{\dfrac{27}{2}} 
=\dfrac{11\times \sqrt{2}\times \sqrt{3}\times 3\sqrt{3}}{3\times 2 \times \sqrt{2}}
=\dfrac{33}{2}
= 16,5$ 
\end{enumerate}

\end{document}
